% GENERATED FILE. Edit canonical lessons, then run npm run generate:books.
% canonical-source-hash: fnv1a64:406f4af78f3c73ec
% canonical-lessons: KA-C181-upayukta, KA-C181-visesa, KA-C181-accumeccina, KA-C181-adhunika, KA-C181-sampradayika

\chapter{Useful, Special, Favourite, Modern, Traditional}
\label{ch:ka-useful-special-favourite-modern-traditional}
\hlchaptermodality{181}

\begin{chapteropening}
\textbf{By the end of this chapter:} \emph{I can say useful, special, favourite, modern and traditional.}

Complete the last lesson of chapter 181: I can say useful, special, favourite, modern and traditional.
\end{chapteropening}

\section[\texorpdfstring{\kn{ಉಪಯುಕ್ತ} (upayukta)}{ಉಪಯುಕ್ತ (upayukta)}]{\kn{ಉಪಯುಕ್ತ} (upayukta) — useful}

\label{lesson:KA-C181-upayukta}



\begin{quote}
Before the new one: say the Kannada for local, then the Kannada for compulsory.
\end{quote}

\subsection*{You'll want to know: \kn{ಉಪಯುಕ್ತ}}
\textbf{\kn{ಉಪಯುಕ್ತ}} — \emph{upayukta} — \textquotedblleft{}useful\textquotedblright{}.

\begin{grammarlens}[title={What you've built}]
One more word for this chapter.
\end{grammarlens}

\subsection*{Your turn}
\begin{itemize}
\raggedright
  \item \emph{Say it:} \emph{upayukta}
  \item \emph{Say it:} \emph{upayukta}, once more
  \item \emph{Say it:} say \emph{upayukta}
  \item \emph{Recall:} say the Kannada for local, then the Kannada for compulsory, then say \emph{upayukta} again
\end{itemize}

\begin{tcolorbox}[breakable,colback=teal!4,colframe=teal!35!black,title={Before you move on}]
What does \kn{ಉಪಯುಕ್ತ} mean? (\textquotedblleft{}Useful\textquotedblright{}.) Say it once more.
\end{tcolorbox}
\section[\texorpdfstring{\kn{ವಿಶೇಷ} (viśēṣa)}{ವಿಶೇಷ (viśēṣa)}]{\kn{ವಿಶೇಷ} (viśēṣa) — special}

\label{lesson:KA-C181-visesa}



\begin{quote}
Before the new one: say the Kannada for compulsory, then the Kannada for useful.
\end{quote}

\subsection*{You'll want to know: \kn{ವಿಶೇಷ}}
\textbf{\kn{ವಿಶೇಷ}} — \emph{viśēṣa} — \textquotedblleft{}special\textquotedblright{}.

\begin{grammarlens}[title={What you've built}]
One more word for this chapter.
\end{grammarlens}

\subsection*{Your turn}
\begin{itemize}
\raggedright
  \item \emph{Say it:} \emph{viśēṣa}
  \item \emph{Say it:} \emph{viśēṣa}, once more
  \item \emph{Say it:} say \emph{viśēṣa}
  \item \emph{Recall:} say the Kannada for compulsory, then the Kannada for useful, then say \emph{viśēṣa} again
\end{itemize}

\begin{tcolorbox}[breakable,colback=teal!4,colframe=teal!35!black,title={Before you move on}]
What does \kn{ವಿಶೇಷ} mean? (\textquotedblleft{}Special\textquotedblright{}.) Say it once more.
\end{tcolorbox}
\section[\texorpdfstring{\kn{ಅಚ್ಚುಮೆಚ್ಚಿನ} (accumeccina)}{ಅಚ್ಚುಮೆಚ್ಚಿನ (accumeccina)}]{\kn{ಅಚ್ಚುಮೆಚ್ಚಿನ} (accumeccina) — favourite}

\label{lesson:KA-C181-accumeccina}



\begin{quote}
Before the new one: say the Kannada for useful, then the Kannada for special.
\end{quote}

\subsection*{You'll want to know: \kn{ಅಚ್ಚುಮೆಚ್ಚಿನ}}
\textbf{\kn{ಅಚ್ಚುಮೆಚ್ಚಿನ}} — \emph{accumeccina} — \textquotedblleft{}favourite\textquotedblright{}.

\begin{grammarlens}[title={What you've built}]
One more word for this chapter.
\end{grammarlens}

\subsection*{Your turn}
\begin{itemize}
\raggedright
  \item \emph{Say it:} \emph{accumeccina}
  \item \emph{Say it:} \emph{accumeccina}, once more
  \item \emph{Say it:} say \emph{accumeccina}
  \item \emph{Recall:} say the Kannada for useful, then the Kannada for special, then say \emph{accumeccina} again
\end{itemize}

\begin{tcolorbox}[breakable,colback=teal!4,colframe=teal!35!black,title={Before you move on}]
What does \kn{ಅಚ್ಚುಮೆಚ್ಚಿನ} mean? (\textquotedblleft{}Favourite\textquotedblright{}.) Say it once more.
\end{tcolorbox}
\section[\texorpdfstring{\kn{ಆಧುನಿಕ} (ādhunika)}{ಆಧುನಿಕ (ādhunika)}]{\kn{ಆಧುನಿಕ} (ādhunika) — modern}

\label{lesson:KA-C181-adhunika}



\begin{quote}
Before the new one: say the Kannada for special, then the Kannada for favourite.
\end{quote}

\subsection*{You'll want to know: \kn{ಆಧುನಿಕ}}
\textbf{\kn{ಆಧುನಿಕ}} — \emph{ādhunika} — \textquotedblleft{}modern\textquotedblright{}.

\begin{grammarlens}[title={What you've built}]
One more word for this chapter.
\end{grammarlens}

\subsection*{Your turn}
\begin{itemize}
\raggedright
  \item \emph{Say it:} \emph{ādhunika}
  \item \emph{Say it:} \emph{ādhunika}, once more
  \item \emph{Say it:} say \emph{ādhunika}
  \item \emph{Recall:} say the Kannada for special, then the Kannada for favourite, then say \emph{ādhunika} again
\end{itemize}

\begin{tcolorbox}[breakable,colback=teal!4,colframe=teal!35!black,title={Before you move on}]
What does \kn{ಆಧುನಿಕ} mean? (\textquotedblleft{}Modern\textquotedblright{}.) Say it once more.
\end{tcolorbox}
\section[\texorpdfstring{\kn{ಸಾಂಪ್ರದಾಯಿಕ} (sāmpradāyika)}{ಸಾಂಪ್ರದಾಯಿಕ (sāmpradāyika)}]{\kn{ಸಾಂಪ್ರದಾಯಿಕ} (sāmpradāyika) — traditional}

\label{lesson:KA-C181-sampradayika}



\begin{quote}
Before the new one: say the Kannada for favourite, then the Kannada for modern.
\end{quote}

\subsection*{You'll want to know: \kn{ಸಾಂಪ್ರದಾಯಿಕ}}
\textbf{\kn{ಸಾಂಪ್ರದಾಯಿಕ}} — \emph{sāmpradāyika} — \textquotedblleft{}traditional\textquotedblright{}.

\begin{grammarlens}[title={What you've built}]
One more word for this chapter.
\end{grammarlens}

\subsection*{Your turn}
\begin{itemize}
\raggedright
  \item \emph{Say it:} \emph{sāmpradāyika}
  \item \emph{Say it:} \emph{sāmpradāyika}, once more
  \item \emph{Say it:} say \emph{sāmpradāyika}
  \item \emph{Recall:} say the Kannada for favourite, then the Kannada for modern, then say \emph{sāmpradāyika} again
\end{itemize}

\begin{tcolorbox}[breakable,colback=teal!4,colframe=teal!35!black,title={Before you move on}]
What does \kn{ಸಾಂಪ್ರದಾಯಿಕ} mean? (\textquotedblleft{}Traditional\textquotedblright{}.) Say it once more.
\end{tcolorbox}
